\scanpage{84}
\mspage{(41)}
\deleted{\[
\tag{112}
\diagram{\begin{tikzcd}[ampersand replacement=\&, column sep=large,row sep=large]
\Pi M^*_{G,G',G''} \arrow[rr,"\gamma_{G,G',G''}"] \arrow[dr] \arrow[dd] \&\& \Pi M_{G''} \arrow[dr,Rightarrow] \\
\& \Pi M^*_{G,G''} \arrow[rr,"\gamma_{G,G''}"] \arrow[dd] \&\& \Pi M_{G''}\\
\Pi M_{G\deleted{,\unsure{G'}}} \arrow[dr,Rightarrow,"\mathrm{can}"'] \\
\& \Pi M_{G,G''}
\end{tikzcd}}
\]
où le carré vertical est 2-cartésien. En d'autres termes, le produit 2-cartésien $\Pi M^*_{G,G',G''}$ défini formellement via (105) comme composé de deux extensions successives de $\Pi M_{G,G'}$, s'interprète canoniquement comme l'extension image inverse de $\Pi M^*_{G,G''}$ sur $\Pi M_{G,G''}$ par
\[
\diagram{\begin{tikzcd}[ampersand replacement=\&]
\Pi M_{G,G'} \arrow[r] \& \Pi M_{G,G''}.
\end{tikzcd}}
\]}

Considérons maintenant le cas d'un morphisme de générisation $G\to G'$, où $G'$ est le graphe de MD \uline{discret} défini par $G$ --- c'est donc la \deleted{\unsure{spécial}}\inserted{générisation} « la plus grossière » de toutes. On \deleted{désignera}\inserted{posera} \deleted{par $S\Pi\unsure{\mathfrak G|\Gamma}M_G$}
\margin{83}
